\section{Conclusion}
\label{sec:conclusion}

This paper proposes frequent losers---firms that never win yet
participate in abnormally many public tenders---as a screening marker for
cover bidding in procurement auctions. Using 4.5 million tender-items
from Brazil's BEC platform (2009--2019), we test two hypotheses of
increasing strength.

\paragraph{H1: FL as screening marker.}
The evidence strongly supports H1: FL presence is a reliable empirical % V7-T2: range with cross-fit
marker for procurement anomalies. FL-present tenders exhibit
3.6--6.6\% higher negotiated prices (3.6\% under temporal
cross-fitting; 6.6\% under full-sample OLS with item, year, and PBU
fixed effects), robust to matching, sensitivity analysis, alternative
FL definitions, and sample restrictions. External validation against
CADE cartel convictions shows that FL firms are 3.5 times more likely
than comparable always-losers to co-participate with known cartelists
($p < 0.001$, permutation test), and the FL--price association
persists after excluding all CADE-involved markets. The FL screen
identifies tenders and firms with characteristics consistent with
collusion.

\paragraph{H2: FL as cover bidders.}
We provide substantial evidence for the causal mechanism. A leave-one-out
instrumental variable strategy yields a price markup of 21\%
(first-stage $F = 396$), exceeding OLS by a factor consistent with
measurement-error attenuation in the binary treatment indicator.
Co-bidding network analysis reveals that the FL--price effect
concentrates in competitive markets---where cover bidding is most
profitable---rather than in concentrated markets where dominant firms
already control pricing. Bajari--Ye tests reject exchangeability ($D = 0.15$, $p < 0.001$) % V7-T2: tender-FE caveat replaced with resolution
and conditional independence (FL pairwise product $= 4.28$,
$t = 81$) for FL bid residuals. Under tender fixed effects, the
FL--non-FL reversal is consistent with Regime~2: cover bids
calibrated near a focal point have low residual covariance after
absorbing tender means, while genuine bidders retain correlated
input-market cost draws. The structural estimation selects
Regime~2 (coordinated cover bidding), with cover-bid dispersion
28\% below genuine-bid dispersion
($\hat{\sigma}_c / \hat{\sigma}_g = 0.72$).
While we cannot rule out all alternative explanations, the convergence
of IV, network, and bid-level evidence makes cover bidding the most
coherent interpretation of the data.

\paragraph{Caveats and limitations.}
The staggered DiD exercise yields imprecise results, reflecting the
difficulty of identifying dynamic treatment effects in markets with
infrequent FL entry and exit. The FL definition is based on a
statistical threshold (median + 1.5 $\times$ IQR) that necessarily
involves some misclassification---though robustness to alternative
thresholds, low-win-rate variants, and cross-fitting suggests this is
not driving the results. As with any cartel screen, strategic adaptation
is a concern: if FL screening were widely adopted, cartels could rotate
cover bidders or allow them occasional wins to avoid detection.

\paragraph{Policy implications.}
The FL screen has three operational advantages for procurement agencies,
particularly in developing countries: (i) it requires only participation
and outcome data, not bid values; (ii) it produces firm-level flags,
enabling investigation of all tenders involving flagged firms; (iii) it
can be combined with bid-level screens in a two-stage workflow---FL
screening first, followed by targeted Bajari--Ye or variance-based
analysis of flagged tenders. The estimated welfare loss ranges from % V7-T2: cross-fit welfare prominent
R\$400 million (cross-fit conservative estimate) to R\$734 million
(OLS) to R\$2.4 billion (IV upper bound) across 2009--2019,
underscoring the potential return on investment from systematic
screening.

\paragraph{Implementation blueprint.}
We outline a practical adoption pathway for procurement agencies.
\emph{Step~1---Flag}: compute the always-loser pool and apply the IQR
threshold to identify FL firms; this requires only participation and
outcome data. \emph{Step~2---Triage}: rank flagged tenders by FL
intensity (number of FL bidders) and co-bidding network concentration
(winner HHI, repeat partners); focus investigative resources on
high-intensity, low-concentration markets where the FL--price effect is
largest. \emph{Step~3---Investigate}: for triaged tenders, apply
bid-level screens (Bajari--Ye, variance tests) to assess coordination
patterns. \emph{Step~4---Monitor}: track FL firms across time to detect
strategic adaptation (e.g., rotation of cover bidders, occasional
wins). This two-stage workflow---firm-level screening followed by
targeted bid-level analysis---reduces the data requirements and
computational burden relative to universal bid-level screening.
For competition authorities, the FL screen can be integrated with
leniency programs by prioritizing markets where FL-based anomalies
coincide with self-reported conduct---allowing limited investigative
resources to be targeted at the highest-probability cases. Importantly,
FL flags are a starting point for investigation, not legal evidence of
collusion. They are best used to trigger further audit and bid-level
forensic analysis rather than as standalone grounds for enforcement
action.

\paragraph{Broader implications.}
This paper complements tender-level forensic screens with a firm-level
criterion operable by any competition authority with access to electronic
procurement records. The convergence of our empirical strategies suggests
cover bidding is the dominant mechanism---a finding with direct
implications for the design of detection and leniency programs.

\paragraph{Future research.}
Four directions merit investigation: (i) external validation on
procurement platforms beyond BEC; (ii) deeper network analysis exploring
shared legal representatives, addresses, and financial links among FL
firms and their frequent co-winners; (iii) dynamic models of how FL
screens interact with leniency programs and cartel stability; (iv)
field experiments where procurement agencies adopt FL screening and
monitor outcomes.
